\section*{अध्याय \ref{ch:b3:green-industrial} --- हरित रसायन और औद्योगिक प्रक्रम}

\begin{solution}{exo:b3:green-industrial:1}
डील्स--ऐल्डर: 100\,\% (एक योग)। एस्टरीकरण: $88.0/(60.0 + 46.0) = 83$\,\%। विटिग:
$104.0/(106.0 + 276.0) = 27$\,\%, द्रव्यमान का अधिकांश भाग ट्राइफ़ेनिलफ़ॉस्फ़ीन ऑक्साइड ले जाता है।
\end{solution}

\begin{solution}{exo:b3:green-industrial:2}
$\mathrm{PMI} = 120/8.0 = 15$; $E = \mathrm{PMI} - 1 = 14$।
\end{solution}

\begin{solution}{exo:b3:green-industrial:3}
किसी दिए गए ताप और सफ़ाई-परिणाम पर एक मानक भार की एक धुलाई। डिटर्जेंटों की
मात्राएँ भिन्न होती हैं: सांद्र पाउडर प्रति धुलाई कम लगता है, अतः समान द्रव्यमान
समान सेवा नहीं देते।
\end{solution}

\begin{solution}{exo:b3:green-industrial:4}
डाइक्लोरोमेथेन के लिए 2-मेथिलटेट्राहाइड्रोफ़्यूरान या एथिल ऐसीटेट; हेक्सेन के लिए हेप्टेन (एथिल
ऐसीटेट के साथ); DMF के लिए, अभिकर्मकों के अनुसार, एथिल ऐसीटेट, 2-मेथिलटेट्राहाइड्रोफ़्यूरान या
डाइमेथिल कार्बोनेट, या \omterm{def:b3:colloids:surfactant}{पृष्ठसक्रियकों} के साथ जल में युग्मन।
\end{solution}

\begin{solution}{exo:b3:green-industrial:5}
1.000 mol अम्ल; 0.800 mol एस्टर (\qty{88.0}{g/mol}): लब्धि 80\,\%। AE 83.0\,\%। SF $= 152/106 = 1.43$।
$\mathrm{RME} = 70.4/152 = 0.463$; $\mathrm{AE} \times \text{लब्धि}/\mathrm{SF} = 0.830 \times 0.80/1.43 = 0.463$।
\end{solution}

\begin{solution}{exo:b3:green-industrial:6}
क्लोरोहाइड्रिन: $44.0/(28.0 + 71.0 + 74.1) = 25$\,\%; प्रत्यक्ष ऑक्सीकरण: 100\,\%। प्रति मोल एथिलीन ऑक्साइड
एक मोल $\ce{CaCl2}$: $\qty{1.0e6}{g}/\qty{44.0}{g/mol} \times \qty{111.1}{g/mol} = \qty{2.5}{t}$ प्रति टन।
\end{solution}

\begin{solution}{exo:b3:green-industrial:7}
$3/8 \times \qty{5.88e4}{mol} \times \qty{44.0}{g/mol} = \qty{0.97}{t}$ $\ce{CO2}$ प्रति टन अमोनिया। हाइड्रोजन के लिए, $\ce{CH4 + 2H2O -> CO2 + 4H2}$:
प्रति मोल $1/4$ मोल, $\qty{5.0e5}{mol}/4 \times \qty{44.0}{g/mol} = \qty{5.5}{t}$ $\ce{CO2}$ प्रति टन हाइड्रोजन।
\end{solution}

\begin{solution}{exo:b3:green-industrial:8}
% ledger: dfg:H2O_l, dfh:H2O-l
एक किलोग्राम \qty{500}{mol} है। $\Delta_rG^\circ = \qty{237.1}{kJ/mol}$ से: $\qty{1.19e5}{kJ} = \qty{32.9}{kWh}$। $\Delta_rH^\circ = \qty{285.8}{kJ/mol}$ से: \qty{39.7}{kWh}
(अंतर वह ऊष्मा है जो परिवेश दे सकता है)।
\end{solution}

\begin{solution}{exo:b3:green-industrial:9}
$\qty{1.0e6}{g}/\qty{146.0}{g/mol} = \qty{6.85e3}{mol}$ $\ce{N2O}$, $\qty{0.30}{t}$; $\times 273 = \qty{82}{t}$ $\ce{CO2}$ तुल्यांक प्रति टन ऐडिपिक अम्ल।
\end{solution}

\begin{solution}{exo:b3:green-industrial:10}
किसी चरण का RME अधिक अभिकर्मक को उपभुक्त गिनता है; यदि अधिक भाग पुनःप्राप्त करके
वापस डाला जाए, तो केवल वास्तव में खोया भाग उपभुक्त होता है, और समग्र द्रव्यमान
दक्षता चरण-मानों के गुणनफल से अधिक होती है। \omterm{def:b3:green-industrial:e-factor}{E-गुणक} में
पुनश्चक्रित धारा अपशिष्ट नहीं है, बशर्ते पुनश्चक्रण परिसीमा के भीतर हो;
तब उसकी ऊर्जा और हानियाँ गिननी होती हैं।
\end{solution}

\begin{solution}{exo:b3:green-industrial:11}
पहले: $20 \times 0.80 = \qty{16}{kg}$ विलायक, 10\,\% हानि: \qty{1.6}{kg} अपशिष्ट। बाद में: \qty{4.0}{kg}, \qty{0.40}{kg} हानि।
पूर्ण \omterm{def:b3:green-industrial:e-factor}{E-गुणक} प्रति किलोग्राम उत्पाद 1.2 kg घटता है (और आसवन
ऊर्जा तीन-चौथाई)।
\end{solution}

\begin{solution}{exo:b3:green-industrial:12}
प्रति \omterm{def:b3:green-industrial:lca}{क्रियात्मक इकाई} और समान परिसीमा के साथ, यह एक अंक के बजाय प्रत्येक प्रभाव-श्रेणी
(जलवायु, भूमि-उपयोग, जल-उपयोग\dots) को अलग-अलग प्रस्तुत करेगा।
निर्णय के लिए इन श्रेणियों को दिए गए भार, जल और भूमि की स्थानीय कमी,
आँकड़ों की अनिश्चितता, और यह जानना चाहिए कि क्या जैव-आधारित कच्चा माल
भोजन से प्रतिस्पर्धा करता है।
\end{solution}

\begin{solution}{pb:b3:green-industrial:1}
\textbf{1.} $\ce{C10H14 + C4H6O3 -> C12H16O + C2H4O2}$ (HF); $\ce{C12H16O + H2 -> C12H18O}$ (रैने निकेल);
$\ce{C12H18O + CO -> C13H18O2}$ (पैलेडियम)।

\textbf{2.} $\ce{C10H14 + C4H6O3 + H2 + CO -> C13H18O2 + C2H4O2}$।

\textbf{3.} अभिकारकों के लिए $134.0 + 102.0 + 2.0 + 28.0 = \qty{266.0}{g/mol}$; आइबुप्रोफ़ेन \qty{206.0}{g/mol}: AE $= 206.0/266.0 = 77$\,\%।

\textbf{4.} आइसोब्यूटिलबेंज़ीन, ऐसीटिक ऐनहाइड्राइड, एथिल क्लोरोऐसीटेट, सोडियम एथॉक्साइड,
जलीय अम्ल ($\ce{H3O+}$), हाइड्रॉक्सिलऐमीन और जल।

\textbf{5.} $134.0 + 102.0 + 122.5 + 68.0 + 19.0 + 33.0 + 2 \times 18.0 = \qty{514.5}{g/mol}$: AE $= 206.0/514.5 = 40$\,\%।

\textbf{6.} ऐसीटिक अम्ल, सोडियम क्लोराइड, एथेनॉल, कार्बन डाइऑक्साइड, जल, अमोनिया (एक
अमोनियम लवण के रूप में), और रससमीकरणमितीय ऐलुमिनियम क्लोराइड से ऐलुमिनियम लवण।

\textbf{7.} $0.71 + 0.55 + 0.010 + 0.14 + 0.05 = \qty{1.46}{t}$।

\textbf{8.} PMI 1.46; $E = 0.46$।

\textbf{9.} $0.46 - 0.31 = 0.15$।

\textbf{10.} PMI $= 2.0 + 1.5 = 3.5$; $E = 2.5$।

\textbf{11.} 100\,\% से कम लब्धियाँ, अधिक अभिकर्मक, विलायक और उत्प्रेरक की हानियाँ, और
उपचार-पश्चात् प्रक्रिया के पदार्थ, सभी ऐसा अपशिष्ट जोड़ते हैं जिसकी \omterm{def:b3:green-industrial:atom-economy}{परमाणु मितव्ययिता} उपेक्षा करती है।

\textbf{12.} अपशिष्ट रोकिए, \omterm{def:b3:green-industrial:atom-economy}{परमाणु मितव्ययिता} अधिकतम कीजिए, रससमीकरणमितीय अभिकर्मकों के बजाय
उत्प्रेरक प्रयुक्त कीजिए, अनावश्यक व्युत्पन्नों और चरणों से बचिए, ऊर्जा बचाइए।

\textbf{13.} रससमीकरणमितीय ऐलुमिनियम क्लोराइड का, जो उपचार-पश्चात् प्रक्रिया में जल-अपघटित होकर ऐलुमिनियम
अपशिष्ट बन जाता है; हाइड्रोजन फ़्लुओराइड उत्प्रेरक और विलायक दोनों है और पुनःप्राप्त करके पुनःप्रयुक्त होता है।

\textbf{14.} रैने निकेल (एक विषमांगी हाइड्रोजनन उत्प्रेरक)।

\textbf{15.} एक पैलेडियम फ़ॉस्फ़ीन संकुल: \cref{ch:b3:homogeneous-catalysis} का
\omterm{def:b3:homogeneous-catalysis:carbonylation}{कार्बोनिलन} रसायन।

\textbf{16.} ऐसीटिक ऐनहाइड्राइड एक ऐसीटिल समूह स्थानांतरित करता है; अणु का दूसरा आधा भाग
ऐसीटिक अम्ल के रूप में निकल जाता है।

\textbf{17.} हाइड्रोजन फ़्लुओराइड अत्यंत विषैला और संक्षारक है, कार्बन मोनोऑक्साइड विषैली और
ज्वलनशील: प्रतिरोधी पदार्थों के बंद उपकरण, संसूचक और प्रशिक्षित
संचालक।

\textbf{18.} प्रत्येक चरण तापन, शीतलन, पृथक्करण और विलायक पुनःप्राप्ति जोड़ता है; तीन
चरणों को इन सबकी कम आवश्यकता होती है।

\textbf{19.} $2.5 \times 3000 = \qty{7500}{t}$ प्रति वर्ष।

\textbf{20.} $0.15 \times 3000 = \qty{450}{t}$ प्रति वर्ष।

\textbf{21.} प्रति वर्ष लगभग \qty{7000}{t} अपशिष्ट बचता है।

\textbf{22.} नहीं: अपशिष्ट की प्रकृति महत्त्वपूर्ण है (एक टन लवणीय जल किसी स्थायी विषैले यौगिक के एक टन
के समान नहीं), और ऊर्जा, उत्सर्जन और अभिकर्मकों के संकट भी।

\textbf{23.} आपूर्ति-शृंखला में पहले होने वाले अभिकर्मकों और ऊर्जा के निर्माण के, वायु और जल में उत्सर्जनों के, और
उत्पाद के उपयोग और निपटान के प्रभाव, प्रति \omterm{def:b3:green-industrial:lca}{क्रियात्मक इकाई}, अनेक
प्रभाव-श्रेणियों में।

\textbf{24.} त्रि-चरणीय उत्प्रेरकी मार्ग की \omterm{def:b3:green-industrial:atom-economy}{परमाणु मितव्ययिता} लगभग 77\,\% है (षट्-चरणीय मार्ग के लिए
40\,\%)।
\end{solution}
